\section*{Hoofdstuk \ref{ch:g10:synthesis-yield} --- Synthese: rendement en zuiverheid}

\begin{solution}{exo:g10:synthesis-yield:1}
$\eta = 4.2 / 5.0 = 0.84$, een \omterm{def:g10:synthesis-yield:yield}{rendement} van \qty{84}{\percent}.
\end{solution}

\begin{solution}{exo:g10:synthesis-yield:2}
De \omterm{def:g7:chemical-reactions:reactant}{beginstoffen} afwegen; \omterm{def:g10:synthesis-yield:reflux}{verwarmen onder terugvloeiing}; \omterm{def:g10:synthesis-yield:vacuum-filtration}{vacuümfiltratie}
van het \omterm{def:g10:synthesis-yield:synthesis}{ruwe product}; \omterm{def:g10:synthesis-yield:recrystallisation}{herkristallisatie}; het smeltpunt meten.
\end{solution}

\begin{solution}{exo:g10:synthesis-yield:3}
Komt het water onderaan binnen, dan vult het de hele mantel voordat het
er bovenaan uit gaat, zodat de koeler over zijn hele lengte gekoeld
wordt. De bovenkant blijft open zodat de opstelling niet gesloten is:
een gesloten opstelling verwarmen zou druk opbouwen en kan haar doen
barsten.
\end{solution}

\begin{solution}{exo:g10:synthesis-yield:4}
Het product moet heel goed \omterm{def:g2:mixing-and-dissolving:dissolve}{oplossen} in het hete \omterm{def:g6:solutions-and-solubility:solute}{oplosmiddel} en nauwelijks
in het koude. De verontreinigingen, die in kleine hoeveelheden aanwezig
zijn, blijven opgelost in het koude \omterm{def:g6:solutions-and-solubility:solute}{oplosmiddel} en komen in het \omterm{def:g3:separating-mixtures:filter}{filtraat}
terecht.
\end{solution}

\begin{solution}{exo:g10:synthesis-yield:5}
Nee: het smelt onder \qty{135}{\celsius} en over een traject van enkele
graden, het teken van een onzuivere vaste stof. Het moet
geherkristalliseerd en gedroogd worden, en het smeltpunt moet opnieuw
gemeten worden.
\end{solution}

\begin{solution}{exo:g10:synthesis-yield:6}
$n = 2.00 / 138.0 = \qty{0.0145}{mol}$ salicylzuur, de \omterm{def:g10:reaction-progress-table:limiting}{beperkende
beginstof}; $m_{\max} = 0.0145 \times 180.0 = \qty{2.61}{g}$;
$\eta = 2.03 / 2.61 = 0.78$, dus \qty{78}{\percent}.
\end{solution}

\begin{solution}{exo:g10:synthesis-yield:7}
De gewogen massa bevat water, dus ze is groter dan de massa aspirine:
het berekende \omterm{def:g10:synthesis-yield:yield}{rendement} is te hoog, hier onmogelijk boven
\qty{100}{\percent}. Het droge \omterm{def:g10:synthesis-yield:synthesis}{ruwe product} bevat nog verontreinigingen
(niet gereageerd salicylzuur, bijproducten), die ook massa toevoegen: het
\omterm{def:g10:synthesis-yield:yield}{rendement} daarvan is ook te hoog vergeleken met het echte \omterm{def:g10:synthesis-yield:yield}{rendement} aan
zuivere aspirine.
\end{solution}

\begin{solution}{exo:g10:synthesis-yield:8}
Ze gaat veel sneller, en ze laat de kristallen veel droger achter, omdat
de zuiging het meeste van de vloeistof wegtrekt; de kristallen zijn ook
makkelijk op de trechter te wassen en van het platte filtreerpapier te
schrapen.
\end{solution}

\begin{solution}{exo:g10:synthesis-yield:9}
Product $3.1/5.0 = 0.62$; aspirine $3.1/5.0 = 0.62$; salicylzuur
$2.2/5.0 = 0.44$. Het product loopt even ver als aspirine en geeft geen
vlek van salicylzuur: het is aspirine, zonder aantoonbaar salicylzuur.
\end{solution}

\begin{solution}{exo:g10:synthesis-yield:10}
\omterm{def:g6:solutions-and-solubility:solute}{Oplosmiddel} A: heet heel goed oplosbaar, koud nauwelijks; \omterm{def:g6:solutions-and-solubility:solute}{oplosmiddel} B
houdt ook koud het grootste deel van het product opgelost. Met A heeft
\qty{3.0}{g} ongeveer $3.0 / 150 = \qty{0.020}{L}$ nodig, dat is
\qty{20}{mL} heet \omterm{def:g6:solutions-and-solubility:solute}{oplosmiddel}; koud blijft er hoogstens
$2 \times 0.020 = \qty{0.04}{g}$ opgelost.
\end{solution}

\begin{solution}{exo:g10:synthesis-yield:11}
Een kokende vloeistof heeft ruimte boven zich nodig: schuim en spatten
mogen de koeler niet bereiken. Kooksteentjes geven de damp kleine
plekjes om belletjes te vormen, zodat de vloeistof rustig kookt in
plaats van met plotselinge uitbarstingen.
\end{solution}

\begin{solution}{exo:g10:synthesis-yield:12}
$0.80 \times 0.70 = 0.56$: \qty{56}{\percent}. Drie stappen van
\qty{90}{\percent}: $0.90^3 = 0.73$, \qty{73}{\percent}. Elke stap kost
product (en tijd en geld), en de verliezen vermenigvuldigen zich: hoe
minder stappen, hoe meer product er aan het eind over is.
\end{solution}

\begin{solution}{exo:g10:synthesis-yield:13}
Hoogstens $3.3 \times 0.050 = \qty{0.17}{g}$ aspirine. IJskoud water lost
nog minder aspirine op, omdat de \omterm{def:g6:solutions-and-solubility:solubility}{oplosbaarheid} daalt met de temperatuur.
\end{solution}

\begin{solution}{exo:g10:synthesis-yield:14}
\begin{enumerate}
  \item $n_0(\text{salicylzuur}) = 2.76 / 138.0 = \qty{0.0200}{mol}$,
        anhydride \qty{0.0500}{mol}: het salicylzuur is beperkend,
        $x_{\max} = \qty{0.0200}{mol}$, en er blijft $0.0500 - 0.0200 =
        \qty{0.0300}{mol}$ anhydride over.
  \item De \omterm{def:g10:synthesis-yield:synthesis}{synthese} geeft \qty{0.0200}{mol} ethaanzuur, de vernietiging
        van de overmaat $2 \times 0.0300 = \qty{0.0600}{mol}$: in totaal
        \qty{0.0800}{mol}.
\end{enumerate}
\end{solution}

\begin{solution}{exo:g10:synthesis-yield:15}
$n(\text{aspirine}) = \num{1.00e6} / 180.0 = \qty{5.56e3}{mol}$. Met een
\omterm{def:g10:synthesis-yield:yield}{rendement} van \qty{85}{\percent} is er aan salicylzuur nodig:
$\num{5.56e3} / 0.85 = \qty{6.54e3}{mol}$, dat is
$\num{6.54e3} \times 138.0 = \qty{9.0e5}{g}$, ongeveer \qty{0.90}{t}.
\end{solution}

\begin{solution}{pb:g10:synthesis-yield:1}
\textbf{1.} C: $7 + 4 = 11$ en $9 + 2 = 11$; H: $6 + 6 = 12$ en
$8 + 4 = 12$; O: $3 + 3 = 6$ en $4 + 2 = 6$. Klopt.

\textbf{2.} Salicylzuur $7 \times 12.0 + 6 \times 1.0 + 3 \times 16.0
= \qty{138.0}{g/mol}$; anhydride $4 \times 12.0 + 6 \times 1.0 +
3 \times 16.0 = \qty{102.0}{g/mol}$; aspirine $9 \times 12.0 + 8 \times
1.0 + 4 \times 16.0 = \qty{180.0}{g/mol}$; ethaanzuur $2 \times 12.0
+ 4 \times 1.0 + 2 \times 16.0 = \qty{60.0}{g/mol}$.

\textbf{3.} \omterm{def:g7:chemical-reactions:reactant}{Beginstoffen}: salicylzuur en ethaanzuuranhydride; gewenst
product: aspirine; bijproduct: ethaanzuur.

\textbf{4.} Verwarmen maakt de reactie sneller; onder terugvloeiing
worden de dampen gecondenseerd en teruggevoerd, zodat er geen \omterm{def:g7:chemical-reactions:reactant}{beginstof}
of product verloren gaat.

\textbf{5.} $n = 3.0 / 138.0 = \qty{0.0217}{mol}$.

\textbf{6.} $m = 6.0 \times 1.08 = \qty{6.48}{g}$;
$n = 6.48 / 102.0 = \qty{0.0635}{mol}$.

\textbf{7.} Salicylzuur: $0.0217 - x$; anhydride $0.0635 - x$;
aspirine en ethaanzuur: $x$. Het salicylzuur is het eerst op:
$x_{\max} = \qty{0.0217}{mol}$, salicylzuur is beperkend.

\textbf{8.} $m_{\max} = 0.0217 \times 180.0 = \qty{3.91}{g}$.

\textbf{9.} Zodat al het salicylzuur in aspirine wordt omgezet: wat er
zou overblijven, blijft met de aspirine gemengd en is moeilijk te
verwijderen, terwijl het overtollige anhydride vernietigd wordt door het
water dat aan het eind wordt toegevoegd, en het ethaanzuur dat daarbij
ontstaat weggewassen wordt.

\textbf{10.} Ze bevat water: \qty{3.6}{g} zou een vals \omterm{def:g10:synthesis-yield:yield}{rendement} van
$3.6 / 3.91 = \qty{92}{\percent}$ geven.

\textbf{11.} $3.3 / 3.91 = 0.84$: \qty{84}{\percent} aan droog \omterm{def:g10:synthesis-yield:synthesis}{ruw
product}, dat niet zuiver is.

\textbf{12.} De verontreinigingen worden verwijderd, en een deel van de
aspirine blijft opgelost in het koude \omterm{def:g6:solutions-and-solubility:solute}{oplosmiddel} en in het waswater.

\textbf{13.} Hoogstens $3.3 \times 0.030 = \qty{0.10}{g}$.

\textbf{14.} $\eta = 2.7 / 3.91 = 0.69$: \qty{69}{\percent}.

\textbf{15.} Een scherp smeltpunt bij de waarde van zuivere aspirine: de
kristallen zijn aspirine, en zuiver.

\textbf{16.} Smelten onder \qty{135}{\celsius} en over een breed traject:
het \omterm{def:g10:synthesis-yield:synthesis}{ruwe product} was onzuiver.

\textbf{17.} Het \omterm{def:g10:synthesis-yield:synthesis}{ruwe product} bevatte aspirine en niet gereageerd
salicylzuur; de \omterm{def:g10:synthesis-yield:recrystallisation}{herkristallisatie} heeft het salicylzuur verwijderd.

\textbf{18.} $R_f = 3.1 / 5.0 = 0.62$.

\textbf{19.} In beide ongeveer evenveel: van \qty{3.91}{g} mogelijk naar
\qty{3.3}{g} \omterm{def:g10:synthesis-yield:synthesis}{ruw product}, en van \qty{3.3}{g} naar \qty{2.7}{g} bij de
\omterm{def:g10:synthesis-yield:recrystallisation}{herkristallisatie}, telkens ongeveer \qty{0.6}{g}. Omdat het \omterm{def:g10:synthesis-yield:synthesis}{ruwe product}
geen zuivere aspirine was, ging er bij het eerste deel van het werk in
werkelijkheid iets meer dan \qty{0.6}{g} aspirine verloren.

\textbf{20.} Het \omterm{def:g10:synthesis-yield:yield}{rendement} van de \omterm{def:g10:synthesis-yield:synthesis}{synthese} is ongeveer \qty{69}{\percent}.
\end{solution}
